\section{Sharp countermodels and scalar boundary}
\label{sec:countermodels}

\begin{example}[Arbitrary tree phases still saturate]
Put arbitrary unit phases and positive weights on the edges of any
tree.  \Cref{thm:tree} gives
\[
 \|H\|=\rho(|H|).
\]
Nonconstant phases alone are not evidence of cancellation.
\end{example}

\begin{example}[The negative sector on a triangle]\label{ex:triangle}
\[
 H=
 \begin{pmatrix}
 0&1&-1\\
 1&0&1\\
 -1&1&0
 \end{pmatrix}
\]
has total holonomy \(-1\) and spectrum \(\{-2,1,1\}\).  It is not
balanced, but it is antibalanced, so
\(\|H\|=2=\rho(|H|)\).
\end{example}

\begin{example}[True frustration on an even cycle]
On \(C_4\), put phase \(-1\) on one edge and \(+1\) on the other
three.  Total holonomy is \(-1\), which fits neither sector because the
cycle length is even.  \Cref{thm:flux} gives
\[
 \|H\|=\sqrt2<2=\rho(|H|).
\]
\end{example}

\begin{example}[Balanced phases with amplitude slack]
Let \(H=(1-\varepsilon)P\), where \(0<\varepsilon<1\) and \(P\) is
connected and nonzero.  Every phase is balanced, but
\[
 \|H\|=(1-\varepsilon)\rho(P)<\rho(P).
\]
This is completion slack, not phase frustration.
\end{example}

\begin{example}[Frustration off the controlling component]
Take the direct sum of a balanced triangle and a frustrated \(C_4\);
both absolute components have Perron radius \(2\).  The triangle still
saturates the global norm.  Finding one frustrated dominant component
does not suffice when another dominant component is balanced.
\end{example}

\begin{example}[Finite strictness can vanish arbitrarily fast]
On a fixed cycle, take total flux \(\Theta_X=X^{-100}\).  Every finite
\(X\) is strictly frustrated in the positive sector, but
\cref{cor:quadratic} gives a defect \(O(X^{-200})\).  Strict inequality
does not imply a fixed power useful to an arithmetic ledger.
\end{example}

\begin{remark}[Why scalar is essential]
For matrix-valued edges, a cycle holonomy such as
\(\diag(1,e^{\mathrm i\vartheta})\ne I\) can still have a common
invariant first coordinate carrying a Perron eigenvector.  Equality may
require a common invariant section rather than trivial holonomy.  For a
general real-linear block \(z\mapsto az+b\overline z\), there is no
single edge phase.  The scalar \(U(1)\) classification must not be
quoted for either setting.
\end{remark}
